\documentclass[11pt]{article}
\pdfinfoomitdate=1
\pdfsuppressptexinfo=15
\pdftrailerid{<74A88D7A8A2BB7EE5665765B699AAA62>
              <74A88D7A8A2BB7EE5665765B699AAA62>}
\usepackage[a4paper,margin=1in]{geometry}
\usepackage{amsmath,amssymb,amsthm,mathtools}
\usepackage[hidelinks]{hyperref}
\usepackage{microtype}

\newtheorem{theorem}{Theorem}
\newtheorem{lemma}[theorem]{Lemma}
\newtheorem{remark}[theorem]{Remark}
\newcommand{\Gal}{\operatorname{Gal}}
\newcommand{\muK}{\mu(K)}
\newcommand{\cchi}{c_\chi}
\newcommand{\Lprime}{L'_{K/k,S}}

\title{Quartic Stark Phase Quantization as a Weak-Solution Equivalence}
\author{Hainan Zhao}
\date{July 31, 2026}

\begin{document}
\maketitle

\begin{abstract}
\noindent\textbf{Claim boundary.}
We isolate two \textsc{proved} consequences of Roblot's quartic
index-solution theorem.  In a cyclic-quartic case where the rank-one
Stark conjecture is already proved, the phase ratio between a Stark
derivative and any Roblot weak solution is a fourth root of unity.  In
an uncertified case satisfying the same algebraic hypotheses, that
phase quantization is equivalent to the full rank-one Stark
conjecture; it is not a finite-precision proof mechanism.  We derive
the root-of-unity factor from the real-place hypothesis and spell out
the even-character, conjugate-character, Fourier-inversion, and
abelian-extension steps.  No census result is claimed here, and this
local draft is not authorized for submission or circulation before
the preregistered counterexample screen exists as its data section.
\end{abstract}

\section{Setup and claim boundary}

Let \(K/k\) be cyclic quartic with
\[
G=\Gal(K/k)=\langle\gamma\rangle,\qquad
\tau=\gamma^2,\qquad \chi(\gamma)=i.
\]
Let \(S\supseteq S(K/k)\) contain a distinguished real place \(v\) of
\(k\) which splits completely in \(K/k\), and fix \(w\mid v\).
We impose Roblot's hypotheses:
\begin{enumerate}
\item[(A1)] \(k\) is totally real, the places of \(K\) above \(v\) are
  real, and the other infinite places of \(K\) are complex;
\item[(A2)] the maximal totally real subfield \(K^+\) has
  \([K:K^+]=2\);
\item[(A3)] every finite prime in \(S\) is ramified or inert in
  \(K/K^+\).
\end{enumerate}
These are exactly the assumptions used in
\cite[pp.~392--393]{Roblot2013}.  We write \(U_K\) for the global
units and use the right Galois exponent \(u^g\).

\begin{lemma}[Real-place root-of-unity rider]\label{lem:roots}
Under \emph{(A1)}, the root-of-unity group of \(K\) is
\[
\muK=\{\pm1\}.
\]
Consequently the denominator \(e=|\muK|\) in the rank-one Stark
formula is \(2\).
\end{lemma}

\begin{proof}
The place \(w\) gives an embedding \(K\hookrightarrow\mathbb R\).
An embedding is injective and sends a root of unity to a real root of
unity.  The only real roots of unity are \(1\) and \(-1\), both of
which lie in \(K\).  Thus \(\muK=\{\pm1\}\).
\end{proof}

For any character \(\psi\) of \(G\) and \(u\in U_K\), put
\[
c_\psi(u)=\frac12\sum_{g\in G}\psi(g)\log|u^g|_w.
\tag{1}\label{eq:coefficient}
\]
For \(a=s\gamma^j\in\{\pm\gamma^j:0\le j<4\}\), where
\(s\in\{\pm1\}\), define signed right exponentiation and its character
value by
\[
u^a=(u^{\gamma^j})^s,\qquad
\chi(a)=s\,i^j.
\]
We separately define the left group-ring action
\[
a\mathbin{\cdot}u:=u^{a^{-1}}.
\tag{2}\label{eq:left-action}
\]
This distinction removes an inverse-convention ambiguity in the
earlier project note.

\begin{lemma}[Pinned covariance]\label{lem:covariance}
For every signed \(a\) and every character \(\psi\),
\[
c_\psi(u^a)=\psi(a)^{-1}c_\psi(u),
\qquad
c_\psi(a\mathbin{\cdot}u)=\psi(a)c_\psi(u).
\tag{3}\label{eq:covariance}
\]
\end{lemma}

\begin{proof}
For \(a=s\gamma^j\), reindexing \(t=r+j\) in
\eqref{eq:coefficient} gives
\[
c_\psi(u^{\gamma^j})
=\frac12\sum_r\psi(\gamma)^r\log|u^{\gamma^{r+j}}|_w
=\psi(\gamma)^{-j}c_\psi(u).
\]
The exponent \(s=-1\) changes every logarithm's sign.  This proves the
first identity; applying it to \(a^{-1}\) proves the second.
\end{proof}

Roblot's Theorem~6.1 \cite[pp.~404--405]{Roblot2013} supplies a class
\(\bar\eta\in\bar U_K^-\) satisfying his index properties (P1) and
(P2), unique under the signed action \(\{\pm\gamma^j\}\).  We fix a
representative \(\eta\) for which
\[
K(\sqrt\eta)/k\quad\text{is abelian}.
\tag{4}\label{eq:weak-abelian}
\]
The same theorem proves
\[
|\Lprime(0,\chi)|=|c_\chi(\eta)|.
\tag{5}\label{eq:weak-absolute}
\]
This is an absolute-value statement, not the oriented Stark identity.

\section{Certified cases}

\begin{theorem}[Certified-case quantization; \textsc{proved}]
\label{thm:certified}
Assume \emph{(A1)--(A3)} and assume the rank-one Stark conjecture has
been proved for \((K/k,S,w)\).  Let \(\epsilon\) be a Stark unit in
the convention
\[
\Lprime(0,\psi)=c_\psi(\epsilon)
\quad\text{for every character }\psi\text{ of }G.
\tag{6}\label{eq:stark}
\]
Then there is a signed element \(h\in\{\pm\gamma^j\}\) such that
\(\bar\eta=h\mathbin{\cdot}\bar\epsilon\), and
\[
\boxed{\frac{\Lprime(0,\chi)}{c_\chi(\eta)}
       =\chi(h)^{-1}\in\mu_4.}
\tag{7}\label{eq:quantized}
\]
\end{theorem}

\begin{proof}
The index formulae for a Stark unit give (P1) and (P2)
\cite[Thm.~2.3]{Roblot2013}.  The uniqueness clause in Roblot's
Theorem~6.1 therefore gives
\(\bar\eta=h\mathbin{\cdot}\bar\epsilon\) for a signed \(h\).
Lemma~\ref{lem:covariance} and \eqref{eq:stark} give
\[
c_\chi(\eta)
=\chi(h)c_\chi(\epsilon)
=\chi(h)\Lprime(0,\chi).
\]
Since the signed quartic character values are precisely \(\mu_4\),
rearrangement proves \eqref{eq:quantized}.
\end{proof}

\begin{remark}[Interpretation]
Theorem~\ref{thm:certified} is a formal consequence of a proved Stark
identity and Roblot's uniqueness theorem.  In an already certified
case it is a convention check, not new evidence for algebraicity.
\end{remark}

\section{Equivalence in uncertified cases}

\begin{theorem}[Weak-solution/Stark equivalence; \textsc{proved}]
\label{thm:equivalence}
Assume \emph{(A1)--(A3)} and fix all conventions above before either
side is evaluated.  The following are equivalent:
\begin{enumerate}
\item the abelian rank-one Stark conjecture holds for \((K/k,S,w)\);
\item there is a signed \(h\in\{\pm\gamma^j\}\) such that
  \[
  \Lprime(0,\chi)=c_\chi(h\mathbin{\cdot}\eta);
  \tag{8}\label{eq:one-character}
  \]
\item the gauge-invariant orbit condition holds:
  \[
  \frac{\Lprime(0,\chi)}{c_\chi(\eta)}\in\mu_4.
  \tag{9}\label{eq:orbit}
  \]
\end{enumerate}
\end{theorem}

\begin{proof}
The implication \(1\Rightarrow2\) follows from
Theorem~\ref{thm:certified} after inverting the signed element if
necessary.  Lemma~\ref{lem:covariance} gives \(2\Rightarrow3\).

Assume (3).  The values of \(\chi\) on the signed group are all of
\(\mu_4\), so choose \(h\) with
\[
\chi(h)=\frac{\Lprime(0,\chi)}{c_\chi(\eta)}
\]
and set \(\bar\epsilon=h\mathbin{\cdot}\bar\eta\).  Then
\eqref{eq:one-character} holds.  We now verify the three propagation
steps needed for the full Stark conjecture.

\smallskip
\noindent\textbf{Step 1: even-component vanishing on both sides.}
Let \(\nu\) be the trivial or quadratic character.  It is even, so
\(\nu(\tau)=1\).  Since \(\bar\eta\in\bar U_K^-\), pairing \(g\) with
\(g\tau\) in \eqref{eq:coefficient} gives
\[
c_\nu(\eta)=0,
\qquad
c_\nu(\epsilon)=0.
\]
For the analytic side, the one-place signature gives
\[
\Lprime(0,\nu)=0
\]
for both even characters; this is the even-character vanishing used
explicitly in the proof of Roblot's Theorem~6.1
\cite[p.~405]{Roblot2013}.  Thus the Stark identity holds on both even
components.

\smallskip
\noindent\textbf{Step 2: conjugate-pair determination and Fourier
inversion.}
All logarithms in \eqref{eq:coefficient} are real, hence
\[
c_{\bar\chi}(\epsilon)=\overline{c_\chi(\epsilon)}.
\]
Artin \(L\)-functions satisfy
\[
\Lprime(0,\bar\chi)
=\overline{\Lprime(0,\chi)}.
\]
Therefore \eqref{eq:one-character} also gives the identity for
\(\bar\chi\).  We now have the Stark identity for all four characters
of \(C_4\).  Fourier inversion on \(G\) determines each component of
the logarithmic vector
\[
\left(\frac12\log|\epsilon^g|_w\right)_{g\in G},
\]
so the complete logarithmic part of the rank-one Stark conjecture
holds.

\smallskip
\noindent\textbf{Step 3: inheritance of the abelian condition.}
Choose an actual representative
\[
\epsilon=\rho\,\eta^{h^{-1}},
\qquad \rho\in\{\pm1\}.
\tag{10}\label{eq:representative}
\]
Conjugation and inversion preserve the abelian extension
\(K(\sqrt\eta)/k\) from \eqref{eq:weak-abelian}: a \(k\)-abelian
extension is stable under \(G\), and adjoining a square root of
\(\eta^{-1}\) gives the same quadratic extension.  If \(\rho=1\),
this proves \(K(\sqrt\epsilon)/k\) is abelian.  If \(\rho=-1\), then
\[
K(\sqrt\epsilon)
\subseteq K(\sqrt\eta)\,k(i).
\]
The extension \(k(i)/k\) is quadratic abelian because \(k\) is totally
real, and a compositum of abelian extensions is abelian.  Hence
\(K(\sqrt\epsilon)/k\) is abelian in either case.  The representative
is a global unit because \(\eta\in U_K\).

The logarithmic identity and the abelian square-root condition are
exactly the abelian rank-one Stark conjecture with \(e=2\), as proved
in Lemma~\ref{lem:roots}.  Thus \(3\Rightarrow1\).
\end{proof}

\section{Numerical and circulation boundaries}

The five archived controls provide only the following
\textsc{observed} statement: after testing both conjugate character
orientations, each independently constructed weak solution matches
one fourth-root rotation against certified \(L'\)-balls.  The old
oriented labels were selected after inspecting the analytic target and
are withdrawn pending exact Artin transport.  They are not used in
either proof above.

A future rigorous screen can refute \eqref{eq:orbit} if the enclosed
analytic value is separated from all four enclosed rotations of the
weak coefficient.  Ball overlap is only
\textsc{certified numerical consistency}, never proof of equality.
No such census screen is contained in this note.

\medskip
\noindent\textbf{Circulation status.}
This is a local research draft.  It must not be submitted, posted, or
circulated before the preregistered B3 screen exists as its data
section.  All email, messaging, submission, and other outbound actions
are human-only.

\section*{AI-assistance disclosure}

AI systems assisted with proof auditing, convention checks, code, and
editing.  The author reviewed the mathematical argument, checked the
cited hypotheses, and is responsible for the claims.

\begin{thebibliography}{9}

\bibitem{Roblot2013}
X.-F.~Roblot,
\emph{Index formulae for Stark units and their solutions},
Pacific J. Math. \textbf{266} (2013), no.~2, 391--417,
\href{https://doi.org/10.2140/pjm.2013.266.391}
{doi:10.2140/pjm.2013.266.391}.

\bibitem{Tate1984}
J.~Tate,
\emph{Les conjectures de Stark sur les fonctions \(L\) d'Artin en
\(s=0\)},
Progress in Mathematics 47, Birkh\"auser, Boston, 1984.

\end{thebibliography}

\end{document}
